% ================================================================
%  SESIÓN 11 — Variables urbanas
%  Almería en Cifras y Formas · Matemáticas 1.º ESO
%  IES Al-Ándalus · 2025-2026
% ================================================================
\documentclass[aspectratio=169, 12pt, spanish]{beamer}
\usepackage{almeria-beamer}

\renewcommand{\SessionNumber}{11}
\renewcommand{\SessionTitle}{Variables urbanas}
\renewcommand{\SessionName}{Sesión 11}
\renewcommand{\SessionObjective}{Identificar variables cuantitativas y cualitativas en el contexto urbano de Almería, distinguiendo variables independientes de dependientes}
\renewcommand{\BlockName}{Bloque 3: Datos y gráficas}

\begin{document}

% ---------------------------------------------------------------
%  PORTADA
% ---------------------------------------------------------------
\begin{frame}[plain]
  \titlepage
\end{frame}

% ---------------------------------------------------------------
%  FRAME 1 — ¿Qué es una variable?
% ---------------------------------------------------------------
\begin{frame}{¿Qué es una variable?}
  \begin{definicionbox}{Variable}
    Una \textbf{variable} es una característica o cualidad que puede tomar
    \textbf{diferentes valores} en los distintos elementos de un conjunto.
    Si todos los elementos tuvieran el mismo valor, no sería variable.
  \end{definicionbox}

  \vspace{0.4cm}
  \begin{columns}[T]
    \begin{column}{0.48\textwidth}
      \begin{ejemplobox}
        \small
        \begin{itemize}
          \item \textbf{Temperatura} de cada día de julio en Almería
          \item \textbf{Número de visitantes} al Alcázar cada mes
          \item \textbf{Tipo de monumento} (árabe, renacentista, moderno\ldots)
          \item \textbf{Barrio} de residencia de cada alumno
        \end{itemize}
      \end{ejemplobox}
    \end{column}
    \begin{column}{0.48\textwidth}
      \begin{notabox}[Dato curioso]
        \small
        El término \textit{variable} viene del latín \textit{variabilis}:
        \textbf{que puede cambiar}.
        En estadística, cada columna de una tabla de datos
        suele representar una variable distinta.
      \end{notabox}
    \end{column}
  \end{columns}
\end{frame}

% ---------------------------------------------------------------
%  FRAME 2 — Clasificación de variables (árbol TikZ)
% ---------------------------------------------------------------
\begin{frame}{Tipos de variables: árbol de clasificación}
  \begin{center}
  \begin{tikzpicture}[
      every node/.style={font=\small},
      nodo/.style={rounded corners=4pt, draw, align=center,
                   inner sep=4pt, text width=2.8cm},
      nodov/.style={nodo, fill=AlmAzul!15, draw=AlmAzul, text=AlmAzul},
      nodoc/.style={nodo, fill=AlmAmbar!20, draw=AlmAmbar!70!black},
      nodod/.style={nodo, fill=AlmTerra!12, draw=AlmTerra},
      flecha/.style={-Stealth, AlmAzulM, line width=1pt},
      scale=0.92
    ]
    % Raíz
    \node[nodov] (var) at (0,0) {\textbf{VARIABLE}};
    % Hijos nivel 1
    \node[nodoc] (cual) at (-3.5,-1.6) {Cualitativa\\\small(describe una\\cualidad)};
    \node[nodoc] (cuant) at (3.5,-1.6) {Cuantitativa\\\small(valor numérico)};
    % Hijos nivel 2
    \node[nodod] (disc) at (2.0,-3.4) {Discreta\\\small números enteros,\\\small contables};
    \node[nodod] (cont) at (5.0,-3.4) {Continua\\\small cualquier valor\\\small en un intervalo};
    % Flechas
    \draw[flecha] (var) -- (cual);
    \draw[flecha] (var) -- (cuant);
    \draw[flecha] (cuant) -- (disc);
    \draw[flecha] (cuant) -- (cont);
    % Ejemplos
    \node[font=\tiny, text=AlmGris, align=center] at (-3.5,-2.9)
      {barrio, tipo de\\monumento, medio\\de transporte};
    \node[font=\tiny, text=AlmGris, align=center] at (2.0,-4.5)
      {turistas/día,\\nº de autobuses};
    \node[font=\tiny, text=AlmGris, align=center] at (5.0,-4.5)
      {temperatura,\\altura, distancia};
  \end{tikzpicture}
  \end{center}
\end{frame}

% ---------------------------------------------------------------
%  FRAME 3 — Cualitativa vs Cuantitativa (detalles)
% ---------------------------------------------------------------
\begin{frame}{Cualitativa y cuantitativa: diferencias clave}
  \begin{columns}[T]
    \begin{column}{0.48\textwidth}
      \begin{definicionbox}{Variable cualitativa}
        \small
        Describe una \textbf{cualidad o categoría}.
        \textbf{No es un número} (aunque a veces se codifique con números).\\[4pt]
        \textit{Ejemplos en Almería:}
        \begin{itemize}\setlength\itemsep{2pt}
          \item Barrio del municipio
          \item Estilo del monumento
          \item Medio de transporte usado
          \item Temporada turística
        \end{itemize}
      \end{definicionbox}
    \end{column}
    \begin{column}{0.48\textwidth}
      \begin{definicionbox}{Variable cuantitativa}
        \small
        Toma \textbf{valores numéricos} con los que tiene sentido operar.\\[4pt]
        \begin{itemize}\setlength\itemsep{2pt}
          \item \textbf{Discreta}: valores enteros y contables\\
            \textit{nº de autobuses de línea, visitantes por día}
          \item \textbf{Continua}: cualquier valor real en un rango\\
            \textit{temperatura, velocidad del viento, consumo de agua}
        \end{itemize}
      \end{definicionbox}
    \end{column}
  \end{columns}

  \vspace{0.3cm}
  \begin{notabox}[Truco para distinguirlas]
    \small
    Pregúntate: \textbf{«¿Tiene sentido calcular la media?»}
    Si sí (media de temperatura = 24{,}3\,°C \correcto),
    es cuantitativa.
    Si no («media de barrios» no tiene sentido \incorrecto),
    es cualitativa.
  \end{notabox}
\end{frame}

% ---------------------------------------------------------------
%  FRAME 4 — Variable independiente y dependiente
% ---------------------------------------------------------------
\begin{frame}{Variable independiente y dependiente}
  \begin{teoriabox}
    \begin{itemize}
      \item \textbf{Variable independiente} (\(x\)): la que controlamos o elegimos.
            Causa el cambio. Va en el \textbf{eje horizontal}.
      \item \textbf{Variable dependiente} (\(y\)): la que cambia \textit{como consecuencia} de \(x\).
            Va en el \textbf{eje vertical}.
    \end{itemize}
  \end{teoriabox}

  \vspace{0.3cm}
  \begin{center}
  \begin{tikzpicture}
    \node[rounded corners=5pt, fill=AlmAzulM!15, draw=AlmAzulM,
          minimum width=3.2cm, minimum height=1cm, font=\bfseries\small]
      (xi) at (0,0) {Mes del año\\{\tiny(variable independiente \(x\))}};
    \node[rounded corners=5pt, fill=AlmAmbar!25, draw=AlmAmbar!70!black,
          minimum width=3.6cm, minimum height=1cm, font=\bfseries\small]
      (yi) at (7,0) {Nº de turistas\\{\tiny(variable dependiente \(y\))}};
    \draw[-Stealth, AlmTerra, line width=2pt] (xi.east) -- node[above,font=\small,color=AlmTerra]
      {\textit{determina}} (yi.west);
  \end{tikzpicture}
  \end{center}

  \vspace{0.2cm}
  \begin{columns}[T]
    \begin{column}{0.55\textwidth}
      \begin{ejemplobox}
        \small
        \textbf{Pregunta clave}: ¿Qué \textit{causa} a qué?\\
        \begin{itemize}
          \item «A más temperatura \textrightarrow\ más consumo de agua»\\
            $x$ = temperatura; $y$ = consumo
          \item «A más lluvia \textrightarrow\ menos turistas»\\
            $x$ = precipitación; $y$ = turistas
        \end{itemize}
      \end{ejemplobox}
    \end{column}
    \begin{column}{0.42\textwidth}
      \begin{notabox}[Fórmula mental]
        \small
        \[\text{Cuanto más } x \;\Rightarrow\; \text{más/menos } y\]
        Si tiene sentido, \(x\) es la independiente.
      \end{notabox}
    \end{column}
  \end{columns}
\end{frame}

% ---------------------------------------------------------------
%  FRAME 5 — Tabla de variables urbanas de Almería
% ---------------------------------------------------------------
\begin{frame}{Variables urbanas de Almería: tabla de ejemplos}
  \begin{center}
  \small
  \renewcommand{\arraystretch}{1.25}
  \begin{tabular}{>{\bfseries\color{AlmAzulM}}l l l l}
    \toprule
    \rowcolor{AlmAzul!12}
    \multicolumn{1}{c}{\textbf{Variable}} & \textbf{Tipo} & \textbf{Rol} & \textbf{Se relaciona con\ldots} \\
    \midrule
    Mes del año & Cuant. discreta & $x$ (independiente) & Nº de visitantes ($y$) \\
    \rowcolor{AlmFondo}
    Temperatura media & Cuant. continua & $x$ & Consumo de agua ($y$) \\
    Barrio & Cualitativa & $x$ & Precio medio de vivienda ($y$) \\
    \rowcolor{AlmFondo}
    Hora del día & Cuant. discreta & $x$ & Intensidad de tráfico ($y$) \\
    Estación del año & Cualitativa & $x$ & Tipo de turista ($y$) \\
    \rowcolor{AlmFondo}
    Nº de bares en calle & Cuant. discreta & $y$ & Mes del año ($x$) \\
    \bottomrule
  \end{tabular}
  \end{center}

  \vspace{0.2cm}
  \begin{notabox}[Observa]
    \small
    Una variable puede ser \textbf{independiente en un par}
    y \textbf{dependiente en otro}.
    El rol depende de la \textit{relación} que estudiamos.
  \end{notabox}
\end{frame}

% ---------------------------------------------------------------
%  FRAME 6 — Gráfico de dispersión: turistas por mes
% ---------------------------------------------------------------
\begin{frame}{Turistas en Almería por mes: diagrama de puntos}
  \begin{center}
  \begin{tikzpicture}
    \begin{axis}[
      almeria,
      width=0.82\linewidth, height=0.55\linewidth,
      xlabel={Mes (1 = enero \ldots 12 = diciembre)},
      ylabel={Turistas (miles)},
      title={Turistas mensuales en Almería (datos aproximados)},
      xtick={1,2,3,4,5,6,7,8,9,10,11,12},
      xticklabels={E,F,M,A,M,J,J,A,S,O,N,D},
      ymin=0, ymax=700,
      ytick={0,100,200,300,400,500,600,700},
      yticklabels={0,100,200,300,400,500,600,700},
    ]
      \addplot[
        only marks,
        mark=*, mark size=3.5pt,
        color=AlmAzulM,
      ] coordinates {
        (1,60) (2,65) (3,80) (4,120) (5,150)
        (6,250) (7,500) (8,600) (9,300) (10,200)
        (11,100) (12,80)
      };
      \addplot[
        AlmTerra!70, dashed, line width=1pt, smooth
      ] coordinates {
        (1,60) (2,65) (3,80) (4,120) (5,150)
        (6,250) (7,500) (8,600) (9,300) (10,200)
        (11,100) (12,80)
      };
      \node[font=\tiny\bfseries, color=AlmTerra, above] at (axis cs:8,600)
        {Máximo: agosto};
      \node[font=\tiny\bfseries, color=AlmAzulM, below right] at (axis cs:1,60)
        {Mínimo: enero};
    \end{axis}
  \end{tikzpicture}
  \end{center}
  \vspace{-0.1cm}
  {\small $x$ = mes (variable independiente) \quad $y$ = turistas en miles (variable dependiente)}
\end{frame}

% ---------------------------------------------------------------
%  FRAME 7 — Ejercicio RECORDAR
% ---------------------------------------------------------------
\begin{frame}{Ejercicio: clasificar variables}
  \begin{recordarbox}
    Clasifica cada variable como \textbf{cualitativa},
    \textbf{cuantitativa discreta} o \textbf{cuantitativa continua}:

    \vspace{0.25cm}
    \renewcommand{\arraystretch}{1.3}
    \begin{tabular}{clc}
      \toprule
      \textbf{N.º} & \textbf{Variable} & \textbf{Tipo} \\
      \midrule
      1 & Altura de la Torre de la Catedral de Almería & \textcolor{AlmGris}{\textit{(completa)}} \\
      2 & Tipo de transporte utilizado & \\
      3 & Visitantes diarios al Alcázar & \\
      4 & Velocidad del viento en el Puerto & \\
      5 & Estilo artístico del monumento & \\
      6 & Consumo de agua por hogar (litros) & \\
      \bottomrule
    \end{tabular}
  \end{recordarbox}

  \vspace{0.2cm}
  \begin{notabox}[Pista]
    \small
    Pregunta 1: la altura puede ser 98{,}36\,m → admite decimales → \textbf{continua}.
    Para las demás, aplica el mismo razonamiento.
  \end{notabox}
\end{frame}

% ---------------------------------------------------------------
%  FRAME 8 — Ejercicio COMPRENDER
% ---------------------------------------------------------------
\begin{frame}{Ejercicio: temperatura y ventas de helados}
  \begin{comprenderbox}
    El ayuntamiento recoge datos sobre la \textbf{temperatura diaria}
    y las \textbf{ventas de helados} en el Paseo de Almería.
    Responde:

    \vspace{0.2cm}
    \begin{enumerate}
      \item ¿Cuál es la variable independiente (\(x\))?
            ¿Por qué?
      \item ¿Cuál es la variable dependiente (\(y\))?
      \item Explica con tus palabras la relación entre ambas.
      \item ¿Qué ocurrirá con las ventas de helados si la temperatura
            sube de 25\,°C a 38\,°C?
            ¿Y si baja a 10\,°C?
    \end{enumerate}
  \end{comprenderbox}

  \vspace{0.3cm}
  \begin{columns}[T]
    \begin{column}{0.48\textwidth}
      \begin{notabox}[Para reflexionar]
        \small
        ¿Podría ser que las ventas de helados
        \textit{causaran} el aumento de temperatura?
        ¿Tiene sentido invertir el par?
      \end{notabox}
    \end{column}
    \begin{column}{0.48\textwidth}
      \begin{notabox}[Vocabulario]
        \small
        «A mayor temperatura, mayores ventas» →
        relación \textbf{directa} o \textbf{positiva}.
      \end{notabox}
    \end{column}
  \end{columns}
\end{frame}

% ---------------------------------------------------------------
%  FRAME 9 — Ejercicio APLICAR
% ---------------------------------------------------------------
\begin{frame}{Ejercicio: inventa pares de variables urbanas}
  \begin{aplicarbox}
    Inventa \textbf{dos pares de variables} propias de Almería:

    \vspace{0.2cm}
    \textbf{Par A}: la variable independiente \(x\) es \textit{el tiempo}
    (hora, día, mes o año).\\
    \textbf{Par B}: la variable independiente \(x\) es
    \textit{una magnitud física} (distancia, temperatura, altura\ldots).

    \vspace{0.2cm}
    Para cada par indica:
    \begin{enumerate}
      \item Nombre y clasificación de \(x\)
      \item Nombre y clasificación de \(y\)
      \item Predicción: «Cuanto más \(x\)\ldots»
    \end{enumerate}
  \end{aplicarbox}

  \vspace{0.2cm}
  \begin{ejemplobox}
    \small
    \textbf{Ejemplo de Par A}: $x$ = mes del año (cuantitativa discreta);
    $y$ = litros de agua consumidos (cuantitativa continua).\\
    Predicción: «En verano, a mayor calor, mayor consumo».
  \end{ejemplobox}
\end{frame}

% ---------------------------------------------------------------
%  FRAME 10 — Ejercicio ANALIZAR
% ---------------------------------------------------------------
\begin{frame}{Ejercicio: una variable oculta}
  \begin{analizarbox}
    El siguiente conjunto de datos muestra turistas en Almería
    en el mes de \textbf{agosto} de varios años:

    \vspace{0.2cm}
    \begin{center}
    \small
    \begin{tabular}{cc}
      \toprule
      \textbf{Año (\(x\))} & \textbf{Turistas agosto (\(y\))} \\
      \midrule
      2019 & 620\,000 \\
      2020 & \textcolor{AlmTerra}{\textbf{180\,000}} \\
      2021 & 390\,000 \\
      2022 & 590\,000 \\
      \bottomrule
    \end{tabular}
    \end{center}

    \vspace{0.2cm}
    \begin{enumerate}
      \item Identifica \(x\) e \(y\) y clasifica cada una.
      \item En 2020 hay una \textbf{caída enorme}.
            El año por sí solo no la explica.
            ¿Qué \textit{variable oculta} podría haber causado ese descenso?
      \item ¿Podría haber más de una variable oculta?
            Razona la respuesta.
    \end{enumerate}
  \end{analizarbox}
\end{frame}

% ---------------------------------------------------------------
%  FRAME 11 — Error frecuente
% ---------------------------------------------------------------
\begin{frame}{¡Cuidado con este error!}
  \begin{errorbox}
    \textbf{Confundir el sentido de la causalidad}:

    \vspace{0.3cm}
    \begin{center}
    \begin{tikzpicture}
      \node[rounded corners=4pt, fill=AlmTerra!15, draw=AlmTerra,
            text width=5cm, align=center, font=\small] (mal)
        at (0,0) {«El \textbf{número de turistas}\\
                  hace que \textbf{cambie el mes}»};
      \node[right=0.5cm of mal, font=\Large, color=AlmTerra] {\incorrecto};
    \end{tikzpicture}

    \vspace{0.4cm}
    \begin{tikzpicture}
      \node[rounded corners=4pt, fill=AlmVerde!12, draw=AlmVerde,
            text width=5cm, align=center, font=\small] (bien)
        at (0,0) {«El \textbf{mes del año}\\
                  influye en el \textbf{número de turistas}»};
      \node[right=0.5cm of bien, font=\Large, color=AlmVerde] {\correcto};
    \end{tikzpicture}
    \end{center}

    \vspace{0.3cm}
    \small
    El \textbf{tiempo} (mes, año, hora) siempre es la variable
    \textit{independiente}: el tiempo avanza solo, no lo controla
    ningún dato de la tabla.
  \end{errorbox}
\end{frame}

% ---------------------------------------------------------------
%  FRAME 12 — Evidencia del día
% ---------------------------------------------------------------
\begin{frame}{Evidencia del día}
  \begin{evidenciabox}
    Completa en tu cuaderno la siguiente \textbf{tabla de variables urbanas}
    con al menos \textbf{5 pares}. Para cada uno indica:

    \vspace{0.2cm}
    \renewcommand{\arraystretch}{1.3}
    \begin{tabular}{>{\bfseries}c l l l l}
      \toprule
      \textbf{Par} & \textbf{Variable \(x\)} & \textbf{Tipo de \(x\)} &
      \textbf{Variable \(y\)} & \textbf{Tipo de \(y\)} \\
      \midrule
      1 & Mes del año & C. discreta & Nº visitantes Alcazaba & C. discreta \\
      2 & & & & \\
      3 & & & & \\
      4 & & & & \\
      5 & & & & \\
      \bottomrule
    \end{tabular}

    \vspace{0.3cm}
    \small
    Indica con \textbf{una flecha} la dirección de la dependencia:
    $x \longrightarrow y$.
    Escribe una frase prediciendo la relación («a mayor $x$\ldots»).
  \end{evidenciabox}
\end{frame}

% ---------------------------------------------------------------
%  FRAME 13 — Resumen de sesión
% ---------------------------------------------------------------
\begin{frame}{Resumen: lo que hemos aprendido hoy}
  \begin{resumenbox}
    \begin{columns}[T]
      \begin{column}{0.48\textwidth}
        \textbf{Tipos de variable:}
        \begin{itemize}
          \item \textbf{Cualitativa}: describe categorías
          \item \textbf{Cuantitativa discreta}: valores enteros
          \item \textbf{Cuantitativa continua}: cualquier valor real
        \end{itemize}
      \end{column}
      \begin{column}{0.48\textwidth}
        \textbf{Independiente vs. dependiente:}
        \begin{itemize}
          \item $x$ \textbf{independiente}: la que controla o causa
          \item $y$ \textbf{dependiente}: la que cambia por $x$
          \item Pregunta: ¿qué \textit{causa} a qué?
        \end{itemize}
      \end{column}
    \end{columns}

    \vspace{0.3cm}
    \separador
    \vspace{0.2cm}
    \small
    \textbf{Contexto de Almería:} mes del año ($x$) \textrightarrow\ turistas ($y$);
    temperatura ($x$) \textrightarrow\ consumo de agua ($y$);
    barrio ($x$) \textrightarrow\ precio vivienda ($y$).
  \end{resumenbox}

  \vspace{0.2cm}
  \begin{center}
    \bloomtag{Recordar}\quad
    \bloomtag{Comprender}\quad
    \bloomtag{Aplicar}\quad
    \bloomtag{Analizar}
  \end{center}
\end{frame}

\end{document}
